\documentclass[aspectratio=169]{beamer}

% Tema y colores
\usetheme{Madrid}
\usecolortheme{default}

% Paquetes necesarios
\usepackage[utf8]{inputenc}
\usepackage[spanish]{babel}
\usepackage{tikz}
\usetikzlibrary {arrows.meta,automata,positioning}

% Configuración de TikZ para autómatas
\tikzset{
    auto,
    node distance=2.5cm,
    initial text=,
    >={Stealth},
    every state/.style={draw=black!50,thick,fill=gray!20}
}

% Información del documento
\title{Counter Machine}
\subtitle{A Python implementation with use examples.}
\author{Jose Luis Cervantes Ariza}
\institute{Universidad de Málaga}
\date{4 de diciembre de 2025}

\begin{document}

% Diapositiva de título
\begin{frame}
    \titlepage
\end{frame}

% Diapositiva de contenido
\begin{frame}{Index}
    \tableofcontents
\end{frame}

% Sección 1
\section{Introduction}

\begin{frame}{Introduction: What is a counter machine?}
    \begin{itemize}
        \item A sequence of \textbf{registers} ${R_1,R_2 ... \in \mathbb{N}}$
        \item A sequence of \textbf{states} ${S_0,S_1,S_2...}$ where $S_0$ stands for the \textbf{halting state}. 
        \item An instruction $\mathbf{(i,+,r)}$ which adds 1 to the register $R_i$ and switches to state $S_r$.
        \item An instruction $\mathbf{(i,-,r,t)}$ which subtracts 1 to the register $R_i$ and switches to state $S_r$, if $R_i=0$ then it switches to state $S_t$ without subtracting.
    \end{itemize}
\begin{center}
    \begin{tabular}{|c|c|c|c|c|}
        \hline
        $R_1$ & $R_2$ & $R_3$ & $R_4$ & ... \\
        \hline
    \end{tabular}
\end{center}
\end{frame}

\begin{frame}{Introduction: A computable function.}
    A (possibly partial) function $f:\mathbb{N}^{k}\rightarrow\mathbb{N}$ is said to be \textbf{computable} if there exists
    a program P (a finite sequence of instructions) for a counter machine such that, given the initial state
    $R_1 = n_1,...,R_k = n_k$, whenever the halting state $S_0$ is reached it turns out that $R_1 = f(n_1,...,n_k)$.
    
\end{frame}

% Sección 2

\section{Implementation}

\begin{frame}{Implementation: The Minsky class.}
The \textbf{Minsky class} has two attributes:
    \begin{itemize}
        \item An array of natural numbers \textbf{registers} which will play the role of the register band (np.ndarray(int)).
        \item A natural number \textbf{max-iterations} which will handle infinite loops (int).
    \end{itemize}
and four methods: 
    \begin{itemize}
        \item A constructor \textbf{init} which takes a natural number indicating the length of \textbf{registers}.
        \item A method \textbf{run} which, given a text file, runs the code written in it.
        \item A method \textbf{debug} which, given a text file, runs the code written in it and prints the result on screen step by step.
        \item A (private) method \textbf{read} which converts text files to ready-to-run counter machine programs. 
    \end{itemize}
\end{frame}

\begin{frame}{Implementation: The Editor class.}
The \textbf{Editor class} has five attributes: the number of registers to be used by the counter machine and another four
attributes related to the arranging of windows, scrollable text, file paths...  \\
It has been implemented making use of the standard Python library \textbf{tkinter}.
\end{frame}

% Sección 2
\section{Use examples}

\begin{frame}{Example: Register exchange.}
    \begin{center}
    \begin{tikzpicture}
        \node[state, initial] (s1) {$S_1$};
        \node[state, accepting, right of=s1] (s0) {$S_0$};
        \node[state, below of=s1] (s3) {$S_3$};
        \node[state, below of=s0] (s2) {$S_2$};
        
        \path[->,dashed] (s1) edge[above] (s0);
        \path[->] (s1) edge[right] node{$R_1-1$} (s2)
              (s2) edge[below] node{$R_2+1$} (s3)
              (s3) edge[left] node{$R_3+1$} (s1);
              
    \end{tikzpicture}
    \end{center}
\end{frame}

\begin{frame}{Example: Subtraction of natural numbers.}
    \begin{center}
    \begin{tikzpicture}
        \node[state, initial] (s1) {$S_1$};
        \node[state, right of=s1] (s2) {$S_2$};
        \node[state,accepting, below of=s2] (s0) {$S_0$};
        
        \path[->,dashed] 
              (s1) edge[bend right, above] (s0)
              (s2) edge[right] (s0);
        \path[->]
              (s1) edge[bend right, below] node{$R_1-1$} (s2)
              (s2) edge[bend right, above] node{$R_2-1$} (s1);
    \end{tikzpicture}
    \end{center}
\end{frame}

\begin{frame}{Example: Euclidean remainder (Initialization.)}
    \begin{center}
    \begin{tikzpicture}
        \node[state, initial] (s1) {$S_1$};
        \node[state, right of=s1] (s4) {$S_4$};
        \node[state, below of=s1] (s3) {$S_3$};
        \node[state, below of=s4] (s2) {$S_2$};

        \node[state, accepting, right of=s4] (s7) {$S_7$};
        \node[state, above of=s4] (s5) {$S_5$};
        \node[state, above of=s7] (s6) {$S_6$};
        
        \path[->,dashed] (s1) edge[above] (s4)
                       (s4) edge[above] (s7);
        \path[->]
              (s1) edge[right] node{$R_1-1$} (s2)
              (s2) edge[below] node{$R_3+1$} (s3)
              (s3) edge[left] node{$R_4+1$} (s1)

              (s4) edge[left] node{$R_2-1$} (s5)
              (s5) edge[above] node{$R_5+1$} (s6)
              (s6) edge[right] node{$R_6+1$} (s4)
              ;
              
    \end{tikzpicture}
    \end{center}
\end{frame}

\begin{frame}{Example: Euclidean remainder (Main loop.)}
    \begin{center}
    \begin{tikzpicture}
        \node[state,initial] (s7) {$S_7$};
        \node[state, right of=s7] (s8) {$S_8$};
        \node[state, right of=s8] (s23) {$S_{23}$};
        \node[state, right of=s23] (s24) {$S_{24}$};  
        \node[state, below of=s7] (s9) {$S_9$};
        \node[state, accepting, below of=s9] (s0) {$S_0$};
        \node[state, right of=s9] (s10) {$S_{10}$};
        \node[state, accepting, right of=s10] (s11) {$S_{11}$};

        \path [->,dashed] (s7) edge[below] (s9)
                       (s9) edge[left] node{$n_1 | n_2$} (s0)
                       (s8) edge[below] node{$n_1>n_2$} (s23)
                       (s23) edge[bend left] (s0);
        \path[->] (s7) edge[bend left, above] node{$R_3-1$} (s8)
              (s8) edge[bend left, below] node{$R_5-1$} (s7)
              (s23) edge[bend left, above] node{$R_6-1$} (s24)
              (s24) edge[bend left, below] node{$R_1+1$} (s23)
              (s9) edge[above] node{$R_5-1$} (s10)
              (s9) edge[below] node{$n_2>n_1$} (s10)
              (s10) edge[above] node{$R_5+1$} (s11);

    \end{tikzpicture}
    \end{center}
\end{frame}

\begin{frame}{Examples: Euclidean remainder (Tidy up of $R_3, R_4$.)}
    \begin{center}
    \begin{tikzpicture}
        \node[state, initial] (s11) {$S_{11}$};
        \node[state, right of=s11] (s14) {$S_{14}$};
        \node[state, accepting, right of= s14] (s16) {$S_{16}$};
        \node[state, below of=s11] (s12) {$S_{12}$};
        \node[state, left of=s12] (s13) {$S_{13}$};
        \node[state, above of=s14] (s15) {$S_{15}$};

        \path [->,dashed] (s11) edge[above] (s14)
                       (s14) edge[above] (s16);
        \path[->] (s11) edge[right] node{$R_4-1$} (s12)
              (s12) edge[below] node{$R_7+1$} (s13)
              (s13) edge[left] node{$R_3+1$} (s11)
              (s12) edge[below] node{$R_7+1$} (s13)
              (s14) edge[bend left, left] node{$R_7-1$} (s15)
              (s15) edge[bend left, right] node{$R_4+1$} (s14);
              
    \end{tikzpicture}
    \end{center}
\end{frame}

\begin{frame}{Example: Euclidean remainder (Tidy up of $R_5, R_6$.)}
    \begin{center}
    \begin{tikzpicture}
        \node[state, initial] (s16) {$S_{16}$};
        \node[state, right of=s16] (s18) {$S_{18}$};
        \node[state] (s21) [right of= s18] {$S_{21}$};
        \node[state, accepting] (s7) [right of= s21] {$S_{7}$};
        \node[state, below of=s16] (s17) {$S_{17}$};
        \node[state, above of=s18] (s19) {$S_{19}$};
        \node[state, above of=s21] (s20) {$S_{20}$};
        \node[state, below of=s21] (s22) {$S_{22}$};

        \path [->,dashed] (s16) edge[above] (s18)
                       (s18) edge[above] (s21)
                       (s21) edge[above] (s7);
        \path[->] (s16) edge[bend right, left] node{$R_6-1$} (s17)
              (s17) edge[bend right, right] node{$R_2+1$} (s16)
              (s18) edge[left] node{$R_5-1$} (s19)
              (s19) edge[above] node{$R_6+1$} (s20)
              (s20) edge[right] node{$R_7+1$} (s18)
              (s21) edge[bend right, left] node{$R_7-1$} (s22)
              (s22) edge[bend right, right] node{$R_5+1$} (s21);
              
    \end{tikzpicture}
    \end{center}
\end{frame}

% Diapositiva final
\begin{frame}[plain]
    \begin{center}
        {\Huge Thank you!}
        
        \vspace{1cm}
        
        {\large Questions, suggestions or congratulations are now welcome.}
    \end{center}
\end{frame}

\end{document}
